\chapter{What money records and what it leaves open}\label{ch:money}

\section{The pump has two kinds of balance}
The clinic manager places two sheets on the table. The first records a payment to the pump workshop. The second records water entering and leaving the clinic's tank. On the first sheet, a debit has a corresponding credit. On the second, an increase at one represented place must be explained by a physical source or a boundary flow. Both sheets speak of balance, but they balance different things.

Imagine that the workshop charges one hundred monetary units for a repair. The clinic pays, the workshop records revenue, and the transaction is complete. Suppose the repaired pump then delivers four units of water. The four water units cannot be obtained by subtracting the two monetary balances. They require a physical record. A change in the negotiated repair price would not, by itself, change the measured water delivery.

This distinction is elementary, yet much confusion arises from ignoring it. The word ``value'' may mean a price, a preference, a benefit, a replacement cost or a declared physical evaluation. The same word does not make the underlying objects identical. Before comparing two values, we should ask which question each number answers and which unit each carries.

\section{Money as a social coordination instrument}
Money allows people to make payments without matching every exchange through direct barter. Prices express exchange terms in a common unit. Saving and borrowing connect claims across time. Accounting makes obligations visible, and contracts specify what parties have agreed to provide. These functions can support exceptionally complicated cooperation among people who never meet.

Consider the clinic's refrigerator. Its components may come from many workshops; its installer needs food and housing; the electricity supplier needs maintenance; the clinic receives income from several sources. A monetary arrangement can coordinate these relationships without making each participant exchange directly with every other participant. Recognising that achievement does not require believing that every resulting allocation is fair or physically sustainable.

The useful question is therefore not whether money is ``real'' in the same sense as water. A contract is real as an institutional relationship, even though it is not a liquid in a tank. The question is how institutional claims are connected to material capabilities. A claim to receive water matters because people accept and enforce it; actual delivery also requires water, equipment and work.

Claims can change without an immediate matching transfer of a particular physical stock. The Bank of England's 2014 account of the modern UK monetary system explains that commercial bank lending can create a corresponding deposit; it rejects the idea that banks merely lend out a fixed pool of prior deposits. Its account also describes constraints on lending and the role of monetary policy. This is an institutional explanation, not a claim that lending creates physical resources without limit.\cite{intro-boe}

For our purposes, the consequence is modest but important. The creation or transfer of a monetary claim does not satisfy the conservation equation for water, fuel or medicine. Such an equation must refer to the represented physical resource. Conversely, conserving water does not settle who is entitled to claim payment. Each account needs its own rules.

\section{A price carries information, but not every kind}
Suppose a replacement valve becomes difficult to obtain. A higher price may encourage repair, substitution or additional supply. It may also prevent a clinic with little purchasing power from obtaining the valve. The physical scarcity and the distribution of purchasing power both affect the situation, but they are different facts.

A price need not be meaningless to be incomplete. It can communicate something about current exchange conditions without being a direct measure of biological need, ecological effect or engineering condition. If two buyers offer different amounts for the last valve, the bids reveal willingness and ability to pay under the arrangement. They do not independently measure the consequences of each buyer going without it.

Nor should we imagine that one can read a price backwards into a unique physical history. The same quoted price might reflect transport expense, bargaining, taxation, quality, scarcity, contractual risk or a temporary discount. One number can result from many combinations. An engineer who needs to know how much electricity a route used must measure or otherwise establish that quantity.

This is why a proposed EBU account should be evaluated as an additional declared description. It cannot acquire accuracy merely by disagreeing with a price. If its state variables are poor, its reference unjustified or its observations wrong, it can also mislead. The comparison should be between explicit information structures, not between an idealised new metric and every defect of actual markets.

\section{Effects outside the immediate exchange}
Imagine two workshops in our invented valley. Both sell a repaired component for the same amount. One contains a cleaning fluid and pays for appropriate treatment; the other allows it to enter a shared stream. If the buyer's payment record is the only account consulted, the two transactions may appear equivalent. The stream record distinguishes them.

This hypothetical example illustrates an external effect: the transaction changes conditions experienced by people or places outside the immediate agreement. It does not prove how common such behaviour is, or specify the best remedy. A physical measurement, a liability rule, a standard, a tax, a collective agreement and a change in production practice might play different roles. The remedy needs evidence about the actual setting.

Positive effects can also be incompletely captured. A person who repairs a shared path may improve access for many people who do not pay that person directly. A family member may perform care whose importance is not represented by a sale. The absence of a market payment is therefore not equivalent to an absence of work or social significance.

An EBU model does not automatically see these effects either. If the stream or the path is outside its declared state, a calculated result may miss them. Extending the representation can improve coverage, but increases measurement and modelling demands. It is preferable to state a limited account honestly than to call it comprehensive because its symbol has an ambitious name.

\section{Existing institutions already connect money and physics}
It would be inaccurate to present the contemporary economy as having no controls, no environmental rules or no public purpose. Institutions already attempt to constrain harmful actions, support essential services and attach consequences to measured physical behaviour. Their effectiveness and coverage vary, but their existence matters to any fair comparison.

For example, the European Commission describes the EU Emissions Trading System as a cap-and-trade arrangement. Covered emissions are monitored and reported; operators surrender allowances corresponding to them. The system therefore connects a declared physical quantity to an institutional obligation and a tradable allowance. This brief description concerns the mechanism, not a claim about its total effectiveness or the status of every proposed reform.\cite{intro-ets}

Notice what this example teaches. The connection requires a defined scope, a unit, a measurement practice, a registry and enforcement rules. It is not generated by the general existence of money. A different physical concern requires its own carefully specified connection. Water availability, access to care and habitat condition are not interchangeable merely because they can appear in financial discussions.

EBU must be situated among such existing efforts. It should explain what information it adds, which decisions it could support and which burdens its use would impose. A useful comparison might eventually show that an existing physical account or a simpler rule performs the relevant task as well. That outcome must remain possible if the programme is to be research.

\section{Four accounts for one repair}
Return to the clinic's pump. We can describe the repair through four linked accounts. The monetary account records the agreed charge and payment. The physical account records the changed equipment condition and any declared resource flows. The service account records whether water became available when needed. The institutional account records authorisation, responsibility and access.

These accounts overlap without collapsing. If the repair is paid for but fails, the monetary event may have occurred while the expected service did not. If a volunteer fixes the pump, physical and service changes can occur without a sale. If the pump works but the clinic excludes a patient, engineering success coexists with an access failure. If the tank fills after rain rather than repair, the physical record must identify the different cause.

The fourth case is especially important for later EBU reasoning. An actor's account should not silently claim a change produced by an independently occurring event. That is why we will compare actions at a declared common baseline and record external changes separately. This distinction is an accounting requirement before it becomes a calculus exercise.

\begin{workedbox}{A payment is neither proof nor disproof of a physical contribution}
A workshop receives one hundred monetary units, and the clinic's store rises by four physical units. To connect these observations, identify the performed action and the measurement interval. The payment alone does not prove the increase; the increase alone does not prove that this workshop caused it. A shared event record makes the proposed connection examinable.
\end{workedbox}

\section{Why financial solvency and physical preservation can diverge}
Take a hypothetical store that begins with ten replacement parts. A manager sells eight and receives more money than the parts originally cost. The monetary result may be favourable. Whether the arrangement remains physically capable depends on replenishment, future needs, alternatives and timing. A positive financial balance does not identify those conditions by itself.

The reverse divergence is equally possible. The manager might keep all ten parts available for emergencies and make no sales. Availability is preserved over that interval, but the financial arrangement may fail to cover wages or rent. A physical account does not pay a bill unless institutions explicitly arrange that relationship. The existence of one kind of reserve cannot be substituted for the other without explaining the conversion.

These examples do not imply that enterprises cannot plan or that monetary analysis ignores the future. They show why planning requires information beyond a single current balance. The EBU research question concerns whether a specific physical valuation and action account can make selected consequences more explicit and useful. It does not begin with a proof that every existing form of planning fails.

Time makes the distinction sharper. Deferring maintenance may improve a current financial result while increasing later risk; performing maintenance may consume resources now to preserve future function. A present-state model must be careful here. It may represent equipment condition directly, but it cannot claim all future benefits merely because the story sounds plausible. Forecasts require an additional model and evidence.

\section{What an EBU result could and could not do}
The signed EBU quantity introduced later describes a change in a declared potential after a declared process burden. It can therefore support a statement such as: ``Relative to this representation and baseline, this accepted action reduced the evaluated burden by this amount.'' The statement is precise enough to check, but narrower than ``this action benefited everyone''.

An institution might choose to display that result, use it in planning, require an explanation for certain patterns or connect it to an incentive. Each choice changes the social arrangement and requires justification. The mathematical account alone does not establish legal tender, ownership, settlement rights, taxation or a conversion rate between EBU and money.

This separation also protects the account. If an essential action has an unfavourable physical result, the institution should be able to support the action without falsifying the record. An emergency delivery over a difficult route may remain necessary. Its consequences can be acknowledged while collective resources cover them. A truthful account is most useful when it survives uncomfortable decisions.

\begin{cautionbox}{Do not infer an institution from an identity}
An accounting equality can establish that defined quantities close. It cannot establish that an actor owns the resulting number, that a bank must accept it, or that people without such numbers should lose access to essential services. Those are additional rules with additional responsibilities.
\end{cautionbox}

\section{A change of price without a change of water}
The clinic and workshop renegotiate their agreement. Yesterday's four-unit delivery cost twenty monetary units; today's identical delivery costs fifteen. Assume, for this example only, identical physical baselines, routes, quantities and declared valuation. The monetary comparison changes. The EBU calculation does not change merely because the negotiated price did.

Now reverse the situation. The charge remains twenty, but the physical baseline differs because Vale already has more water. The same delivered quantity may then have a different represented field effect. The monetary record can remain unchanged while the physical evaluation changes. These paired examples show why a fixed conversion between price and EBU cannot be assumed from their common association with a transaction.

A future institution might nevertheless choose an explicit conversion or reward rule. That would be a policy linking two accounts, with consequences for incentives, access and risk. It would need to explain how revisions are made and what happens when the physical result is negative or uncertain. The link could be studied, but it should not be hidden inside the definition of the physical unit. Keeping the accounts distinct allows readers to inspect whether the proposed institutional link is sensible rather than treating it as an arithmetic necessity.

\section{Common misunderstandings}
``Money measures nothing real'' is too crude. Monetary records describe real institutional relationships and can organise real provision. Their physical interpretation depends on what they are connected to.

``A high price proves high physical burden'' also goes too far. The price and the burden may be associated in a setting, but one cannot be deduced from the other without a model of that setting. A donated resource can carry a substantial physical burden; an expensive service can use little material.

``A physical account can replace every monetary function'' is unsupported. Measuring an action does not automatically supply credit, insurance, dispute resolution or a workable payment system. A proposed institutional extension must specify how those functions would operate.

``If the physical account is exact, institutions become objective'' confuses arithmetic with choice. The representation, reference and permitted uses still require reasons. Exact calculation makes the consequences of those declarations clearer; it does not remove the declarations.

\section{Guided problems and answers}
\subsection{Same price, different event}
Two couriers charge the same amount to deliver the same quantity. One uses a direct route and the other a detour. Can the payment record distinguish their physical effects?

The common price does not establish identical effects. A physical comparison requires the declared routes, carriers, quantities and boundaries. It also requires evidence about what happened. If the teaching model represents only the source and destination stock change, both routes may have the same represented endpoint value; the unrepresented route difference must remain a limitation or be represented through an additional justified account.

\subsection{No income, valuable maintenance}
A community repairs a path through unpaid work. Does zero monetary revenue imply zero EBU?

No inference follows in either direction. The absence of revenue describes the payment arrangement. An EBU calculation would require an appropriate state, reference, physical change and process-burden declaration. Calling the repair socially valuable does not supply those data. The right response is to distinguish the ethical judgement, the physical observations and any mathematical valuation.

\subsection{A claim and a tank}
A bank account increases while the town's water stocks remain unchanged. Has conservation been violated?

Not merely for that reason. The monetary and water accounts track different objects. A violation would occur if the water account reported an unexplained physical creation or disappearance under a closed boundary. Whether the monetary change is valid depends on its own institutional rules. The exercise is a reminder to identify the conserved quantity before applying conservation language.

\section{What to carry forward}
Money is a powerful instrument for organising claims, exchange and time. Its records need explicit connections to particular physical and social objectives. EBU investigates one such physical account and must justify its usefulness on its own terms. The next chapter turns to a consequence that a total or a price can easily hide: resources may exist while people still lack effective access to them.
